%%%%%%%%%%%%%%%%%%%%%%%%%%%%%%%%%%%%%%%%%%%%%%%%%%%%%%%%%%%%%%%%%%%%%%%%%%%%%%
%% § Definiciones                                                           %%
%%%%%%%%%%%%%%%%%%%%%%%%%%%%%%%%%%%%%%%%%%%%%%%%%%%%%%%%%%%%%%%%%%%%%%%%%%%%%%

% Código Lean
% ===========

\let\verbatim\leancode
\let\endverbatim\endleancode

%%% 1. Estilo de teoremas, lemas, etc. (texto en cursiva)
\theoremstyle{plain}
\newtheorem{theorem}{Teorema}[section]
\newtheorem{teorema}{Teorema}[section]
\newtheorem{lema}[theorem]{Lema}
\newtheorem{lemma}[theorem]{Lema}
\newtheorem{corolario}[theorem]{Corolario}
\newtheorem{proposicion}[theorem]{Proposición}

%%% 2. Estilo de definiciones y ejemplos (texto normal)
\theoremstyle{definition}
\newtheorem{definicion}[theorem]{Definición}
\newtheorem{ejemplo}[theorem]{Ejemplo}

%%% 3. Estilo de notas (título en cursiva, texto normal)
\theoremstyle{remark}
\newtheorem{nota}[theorem]{Nota}
\newtheorem{claim}[theorem]{Nota}

%%% 4. Personalización de la demostración
\addto\captionsspanish{
  \renewcommand{\proofname}{\textbf{Demostración:}}
}

% ── Macros matemáticos ────────────────────────────────────────────────────────
\newcommand{\N}{\mathbb{N}}
\newcommand{\R}{\mathbb{R}}
\newcommand{\abs}[1]{\left|#1\right|}

% Definir \square para el final de demostraciones
\newcommand{\square}{\qedsymbol}

% ---------- Diálogo ----------
\newcommand{\socrates}[1]{\medskip\noindent\textbf{SÓCRATES:} #1}
\newcommand{\simplicio}[1]{\medskip\noindent\textbf{SIMPLICIO:} #1}

%%% Local Variables:
%%% mode: latex
%%% TeX-master: t
%%% End:
